\chapter{Instalación en servidor físico: Linux}
\label{cap:linux}

Instalación sin Docker en Debian 12 o Ubuntu 24.04 LTS. Cada componente queda como un servicio
de \texttt{systemd}. Todo se instala desde los repositorios oficiales de cada programa.

\begin{xltabular}{\linewidth}{@{}l L@{}}
	\toprule
	\textbf{Ruta} & \textbf{Contenido} \\
	\midrule
	\endhead
	\texttt{/opt/digaipe/app} & Aplicación (propiedad de \texttt{root}, solo lectura para el
	servicio) \\
	\texttt{/opt/digaipe/app/var/uploads} & Archivos de admisiones cargados (escritura del
	servicio) \\
	\texttt{/etc/digaipe/} & Archivos de configuración con contraseñas (solo \texttt{root}) \\
	\texttt{/var/backups/digaipe} & Respaldos (mejor en otro disco montado ahí) \\
	\texttt{/usr/local/sbin/digaipe-*} & Comandos de administración \\
	\bottomrule
\end{xltabular}

\section{Preparar el sistema}

Zona horaria, usuarios de servicio y carpetas:

\begin{codigo}
sudo timedatectl set-timezone America/Caracas
sudo useradd --system --home-dir /opt/digaipe --shell /usr/sbin/nologin digaipe
sudo useradd --system --no-create-home --shell /usr/sbin/nologin --groups digaipe digaipe-respaldo
sudo mkdir -p /opt/digaipe /etc/digaipe /var/backups/digaipe
sudo chmod 700 /etc/digaipe
sudo chown digaipe-respaldo:digaipe-respaldo /var/backups/digaipe
sudo chmod 700 /var/backups/digaipe
\end{codigo}

\begin{itemize}
	\item \texttt{digaipe} ejecuta la aplicación.
	\item \texttt{digaipe-respaldo} ejecuta los respaldos. Pertenece al grupo \texttt{digaipe} para
	      poder \emph{leer} los archivos cargados, pero no puede modificarlos, y la aplicación no
	      puede leer los respaldos.
\end{itemize}

Si los respaldos van en un segundo disco, móntelo en \cod{/var/backups/digaipe} (entrada en
\cod{/etc/fstab}) \emph{antes} de los comandos \texttt{chown} y \texttt{chmod} de arriba.

\section{Instalar Node.js 24}

Los repositorios de Debian y Ubuntu traen versiones de Node.js más antiguas. Use el repositorio
de NodeSource, que se actualiza con \texttt{apt}:

\begin{codigo}
curl -fsSL https://deb.nodesource.com/setup_24.x -o /tmp/nodesource_setup.sh
sudo bash /tmp/nodesource_setup.sh
sudo apt install -y nodejs
\end{codigo}

El script solo agrega el repositorio y su clave; puede revisarlo
(\cod{less /tmp/nodesource_setup.sh}) antes de ejecutarlo.

\begin{comprobar}
	\texttt{node -v} debe mostrar \texttt{v24.<x>}, y \texttt{npx -v}, un número de versión.
\end{comprobar}

\section{Instalar MySQL 8.4}

Debian no trae MySQL (solo MariaDB) y Ubuntu trae MySQL 8.0. Use el repositorio APT oficial de
MySQL. Descargue el paquete de configuración más reciente desde
\texttt{https://dev.mysql.com/downloads/repo/apt/} e instálelo:

\begin{codigo}
sudo apt install -y gnupg wget
wget https://dev.mysql.com/get/mysql-apt-config_<versión>_all.deb -O /tmp/mysql-apt-config.deb
sudo apt install -y /tmp/mysql-apt-config.deb
\end{codigo}

En el menú que aparece, elija \emph{MySQL Server \& Cluster} → \textbf{mysql-8.4-lts} y luego
\emph{Ok}. Si se equivoca, vuelva al menú con \texttt{sudo dpkg-reconfigure mysql-apt-config}.
Después:

\begin{codigo}
sudo apt update
sudo apt install -y mysql-server
\end{codigo}

El instalador pide la contraseña de \texttt{root} de MySQL: use \texttt{MYSQL\_ROOT\_PASSWORD}.
Si pregunta por el método de autenticación, elija el recomendado (\texttt{caching\_sha2\_password}).

Restrinja MySQL al propio servidor creando \cod{/etc/mysql/mysql.conf.d/digaipe.cnf}:

\begin{codigo}[label={/etc/mysql/mysql.conf.d/digaipe.cnf}]
[mysqld]
# Solo conexiones desde este servidor: la aplicación y los respaldos.
bind-address = 127.0.0.1
mysqlx = OFF
\end{codigo}

Reinicie y cree la base de datos:

\begin{codigo}
sudo systemctl restart mysql
sudo systemctl enable mysql
mysql -h 127.0.0.1 -u root -p -e "CREATE DATABASE IF NOT EXISTS digaipe CHARACTER SET utf8mb4 COLLATE utf8mb4_0900_ai_ci"
\end{codigo}

\begin{comprobar}
	\begin{codigo}
mysql -h 127.0.0.1 -u root -p -e "SELECT VERSION(); SHOW DATABASES LIKE 'digaipe'"
sudo ss -ltnp | grep mysqld
	\end{codigo}
	La versión debe empezar por \texttt{8.4}, la base \texttt{digaipe} debe existir y
	\texttt{mysqld} debe escuchar solo en \texttt{127.0.0.1:3306}.
\end{comprobar}

\section{Instalar la aplicación}

Clone el repositorio en \cod{/opt/digaipe/app} (capítulo~\ref{cap:preparacion}) y construya:

\begin{codigo}
cd /opt/digaipe/app
sudo npm ci
sudo npm run build
sudo mkdir -p /opt/digaipe/app/var/uploads
sudo chown -R digaipe:digaipe /opt/digaipe/app/var
sudo chmod 2750 /opt/digaipe/app/var /opt/digaipe/app/var/uploads
\end{codigo}

El bit \texttt{2} (\emph{setgid}) hace que las carpetas nuevas de \cod{var/uploads} hereden el
grupo \texttt{digaipe}, de modo que el usuario de respaldos siempre pueda leerlas.

El código, \texttt{node\_modules} y \texttt{build} quedan propiedad de \texttt{root}: el
servicio puede leerlos pero no modificarlos. Solo \texttt{var} es suyo.

\begin{comprobar}
	\texttt{npm run build} debe terminar sin errores, y \texttt{ls -a /opt/digaipe/app} debe
	mostrar la carpeta \texttt{build}. \textbf{No} debe haber un archivo \texttt{.env}: la
	configuración va en \cod{/etc/digaipe/}.
\end{comprobar}

\begin{nota}
	El paquete incluye las dependencias de desarrollo (\texttt{drizzle-kit}, \texttt{tsx}) a
	propósito: la instalación las usa para aplicar migraciones, crear las cuentas de MySQL y
	verificar la auditoría.
\end{nota}

\section{Archivos de configuración}
\label{sec:linux-config}

La instalación usa tres archivos, uno por cada cuenta de MySQL, para que cada proceso conozca
solo la contraseña que necesita. Créelos con \texttt{sudo nano <archivo>} y reemplace los
valores \texttt{<...>}.

\begin{codigo}[label={/etc/digaipe/setup.env  (instalación, restauración y auditoría: root de MySQL)}]
MYSQL_ROOT_PASSWORD=<MYSQL_ROOT_PASSWORD>
DATABASE_URL=mysql://root:<MYSQL_ROOT_PASSWORD>@127.0.0.1:3306/digaipe
APP_DB_USER=digaipe_app
APP_DB_PASSWORD=<APP_DB_PASSWORD>
BACKUP_DB_USER=digaipe_respaldo
BACKUP_DB_PASSWORD=<BACKUP_DB_PASSWORD>
READONLY_DB_USER=digaipe_lectura
READONLY_DB_PASSWORD=
ORIGIN=https://<DIGAIPE_HOST>
BETTER_AUTH_SECRET=<BETTER_AUTH_SECRET>
BOOTSTRAP_ADMIN_EMAIL=<correo del administrador>
BOOTSTRAP_ADMIN_PASSWORD=<BOOTSTRAP_ADMIN_PASSWORD>
BOOTSTRAP_ADMIN_NAME=Administrador
\end{codigo}

\begin{codigo}[label={/etc/digaipe/digaipe.env  (aplicación)}]
NODE_ENV=production
DATABASE_URL=mysql://digaipe_app:<APP_DB_PASSWORD>@127.0.0.1:3306/digaipe
ORIGIN=https://<DIGAIPE_HOST>
BETTER_AUTH_SECRET=<BETTER_AUTH_SECRET>
# Solo Caddy, en este mismo servidor, puede conectarse a la aplicación.
HOST=127.0.0.1
PORT=3000
# IP real del usuario, tomada de la entrada que añade Caddy (auditoría y límite de intentos).
ADDRESS_HEADER=X-Forwarded-For
XFF_DEPTH=1
# Las cargas de admisiones (hoja de cálculo + PDF) superan el límite por defecto de 512K.
BODY_SIZE_LIMIT=50M
TZ=America/Caracas
\end{codigo}

\begin{codigo}[label={/etc/digaipe/backup.env  (respaldos)}]
DB_HOST=127.0.0.1
DB_NAME=digaipe
BACKUP_DB_USER=digaipe_respaldo
BACKUP_DB_PASSWORD=<BACKUP_DB_PASSWORD>
BACKUP_ROOT=/var/backups/digaipe
UPLOADS_DIR=/opt/digaipe/app/var/uploads
BACKUP_INTERVAL_DAYS=7
BACKUP_HOUR_UTC=6
BACKUP_KEEP_DAYS=70
BACKUP_KEEP_MIN=10
BACKUP_MIN_FREE_MB=512
\end{codigo}

\begin{codigo}
sudo chown root:root /etc/digaipe/setup.env /etc/digaipe/digaipe.env /etc/digaipe/backup.env
sudo chmod 600 /etc/digaipe/setup.env /etc/digaipe/digaipe.env /etc/digaipe/backup.env
\end{codigo}

Los servicios leen estos archivos como \texttt{root} al arrancar (\texttt{systemd}), así que no
hace falta que los usuarios de servicio puedan leerlos.

\begin{aviso}
	\texttt{ORIGIN} y \texttt{BETTER\_AUTH\_SECRET} deben ser idénticos en \texttt{setup.env} y en
	\texttt{digaipe.env}, y \texttt{APP\_DB\_PASSWORD} / \texttt{BACKUP\_DB\_PASSWORD}, idénticas
	en el archivo de instalación y en el que las usa. Si difieren, la aplicación o los respaldos
	no podrán conectarse.
\end{aviso}

\section{Comandos de administración}

Estos tres scripts cargan \texttt{setup.env} y ejecutan, como \texttt{root} de MySQL, la
instalación, la auditoría y la restauración. Hacen lo mismo que los servicios \texttt{setup},
\texttt{audit} y \texttt{restore} del modo Docker.

\begin{codigo}[label={/usr/local/sbin/digaipe-setup}]
#!/bin/bash
# Migraciones, triggers de auditoría, cuentas de MySQL, administrador inicial y datos base.
# Se puede repetir sin riesgo. Obligatorio después de restaurar o actualizar.
set -euo pipefail
cd /opt/digaipe/app
set -a; . /etc/digaipe/setup.env; set +a
exec sh scripts/docker/setup.sh
\end{codigo}

\begin{codigo}[label={/usr/local/sbin/digaipe-audit}]
#!/bin/bash
# Uso: digaipe-audit verify [--anchors <archivo>]
#      digaipe-audit anchor --anchors <archivo> [--note <texto>]
set -euo pipefail
cd /opt/digaipe/app
set -a; . /etc/digaipe/setup.env; set +a
exec npx tsx scripts/audit-ledger.ts "$@"
\end{codigo}

\begin{codigo}[label={/usr/local/sbin/digaipe-restaurar}]
#!/bin/bash
# Uso: digaipe-restaurar <respaldo> <respaldo>   (el nombre dos veces, como confirmación)
set -euo pipefail
set -a; . /etc/digaipe/backup.env; . /etc/digaipe/setup.env; set +a
RESTORE_CONFIRM="${2:-}" bash /opt/digaipe/app/deploy/backup/restore.sh "${1:?Uso: digaipe-restaurar <respaldo> <respaldo>}"
# restore.sh copia los archivos como root; la aplicación necesita poder escribirlos.
chown -R digaipe:digaipe /opt/digaipe/app/var/uploads
\end{codigo}

\begin{codigo}
sudo chmod 700 /usr/local/sbin/digaipe-setup /usr/local/sbin/digaipe-audit /usr/local/sbin/digaipe-restaurar
\end{codigo}

\section{Ejecutar la instalación}

\begin{codigo}
sudo /usr/local/sbin/digaipe-setup
\end{codigo}

\begin{comprobar}
	La salida debe mostrar, en orden y sin errores: \texttt{==> Migraciones},
	\texttt{==> Triggers de auditoría}, \texttt{==> Usuarios de base de datos y permisos}
	(con líneas \texttt{Granted digaipe\_app ...} y \texttt{Granted digaipe\_respaldo ...}),
	\texttt{==> Administrador inicial} y \texttt{==> Datos base}.
\end{comprobar}

\section{Servicio de la aplicación}

\begin{codigo}[label={/etc/systemd/system/digaipe.service}]
[Unit]
Description=DIGAIPE (aplicacion web)
After=network-online.target mysql.service
Wants=network-online.target
Requires=mysql.service

[Service]
Type=simple
User=digaipe
Group=digaipe
WorkingDirectory=/opt/digaipe/app
EnvironmentFile=/etc/digaipe/digaipe.env
ExecStart=/usr/local/bin/node build
Restart=on-failure
RestartSec=5
# Archivos nuevos legibles por el grupo (respaldos), no por otros usuarios.
UMask=0027
NoNewPrivileges=true
PrivateTmp=true
ProtectHome=true
ProtectSystem=strict
ReadWritePaths=/opt/digaipe/app/var

[Install]
WantedBy=multi-user.target
\end{codigo}

\texttt{WorkingDirectory} es obligatorio: la aplicación busca la plantilla PDF en
\cod{static/} y guarda las cargas en \cod{var/uploads/}, ambas relativas a esa carpeta.

\begin{codigo}
sudo systemctl daemon-reload
sudo systemctl enable --now digaipe
\end{codigo}

\begin{comprobar}
	\begin{codigo}
systemctl status digaipe
curl -sI http://127.0.0.1:3000/ | head -n 1
	\end{codigo}
	El servicio debe estar \texttt{active (running)} y \texttt{curl} debe responder con una línea
	\texttt{HTTP/1.1}. Si falla, vea \texttt{journalctl -u digaipe -n 50}.
\end{comprobar}

\section{Caddy (HTTPS)}

Desde el repositorio oficial de Caddy:

\begin{codigo}
sudo apt install -y debian-keyring debian-archive-keyring apt-transport-https curl
curl -1sLf 'https://dl.cloudsmith.io/public/caddy/stable/gpg.key' | sudo gpg --dearmor -o /usr/share/keyrings/caddy-stable-archive-keyring.gpg
curl -1sLf 'https://dl.cloudsmith.io/public/caddy/stable/debian.deb.txt' | sudo tee /etc/apt/sources.list.d/caddy-stable.list
sudo apt update
sudo apt install -y caddy
\end{codigo}

El paquete crea el usuario \texttt{caddy} y el servicio \texttt{caddy}. Reemplace
\cod{/etc/caddy/Caddyfile} por:

\begin{codigo}[label={/etc/caddy/Caddyfile}]
# HTTPS de DIGAIPE: Caddy actúa como autoridad certificadora local (red interna).
{
	local_certs
	# La raíz se instala a mano en las PCs cliente, no en este servidor.
	skip_install_trust
}

# Debe coincidir exactamente con ORIGIN (sin https://).
<DIGAIPE_HOST> {
	tls internal
	reverse_proxy 127.0.0.1:3000
	header -Server
}
\end{codigo}

\begin{codigo}
sudo caddy validate --config /etc/caddy/Caddyfile
sudo systemctl enable caddy
sudo systemctl restart caddy
\end{codigo}

La autoridad certificadora queda en
\cod{/var/lib/caddy/.local/share/caddy/pki/authorities/local/}. \textbf{Esa carpeta debe
estar en las copias de seguridad del servidor}: si se pierde, Caddy crea otra autoridad y hay
que volver a instalar la raíz en todas las PCs (capítulo~\ref{cap:certificado}).

\section{Servicio de respaldos}

\begin{codigo}[label={/etc/systemd/system/digaipe-respaldo.service}]
[Unit]
Description=DIGAIPE (respaldo semanal)
After=mysql.service
Requires=mysql.service

[Service]
Type=simple
User=digaipe-respaldo
Group=digaipe-respaldo
SupplementaryGroups=digaipe
EnvironmentFile=/etc/digaipe/backup.env
ExecStart=/usr/bin/bash /opt/digaipe/app/deploy/backup/worker.sh
Restart=always
RestartSec=60
NoNewPrivileges=true
PrivateTmp=true
ProtectHome=true
ProtectSystem=strict
ReadWritePaths=/var/backups/digaipe

[Install]
WantedBy=multi-user.target
\end{codigo}

\begin{codigo}
sudo systemctl daemon-reload
sudo systemctl enable --now digaipe-respaldo
\end{codigo}

\begin{comprobar}
	\texttt{journalctl -u digaipe-respaldo -n 20} debe mostrar
	\texttt{Backup worker started}, \texttt{Backup due: no backup exists yet} y, al cabo de unos
	segundos, \texttt{Backup finished.} \texttt{sudo ls /var/backups/digaipe} debe mostrar una
	carpeta \texttt{digaipe-<fecha>}.
\end{comprobar}

\section{Cortafuegos}

\begin{codigo}
sudo ufw allow 22/tcp
sudo ufw allow 80/tcp
sudo ufw allow 443/tcp
sudo ufw enable
sudo ufw status
\end{codigo}

Omita la regla 22 si no administra el servidor por SSH. Si el sistema usa \texttt{nftables} o
\texttt{firewalld} en lugar de \texttt{ufw}, aplique las mismas reglas con esa herramienta.

\section{Siguiente paso}

Continúe con el certificado HTTPS (capítulo~\ref{cap:certificado}) y la verificación
(capítulo~\ref{cap:verificacion}).

\section{Operación diaria}

\begin{xltabular}{\linewidth}{@{}l L@{}}
	\toprule
	\textbf{Tarea} & \textbf{Comando} \\
	\midrule
	\endhead
	Ver estado & \texttt{systemctl status digaipe caddy mysql digaipe-respaldo} \\
	Ver registros de la aplicación & \texttt{journalctl -u digaipe -f} \\
	Reaplicar la instalación & \texttt{sudo /usr/local/sbin/digaipe-setup} y luego
	\texttt{sudo systemctl restart digaipe} \\
	Detener & \texttt{sudo systemctl stop digaipe digaipe-respaldo} \\
	Iniciar & \texttt{sudo systemctl start digaipe digaipe-respaldo} \\
	\bottomrule
\end{xltabular}
